\documentclass[border=10pt]{standalone}
\usepackage{tikz,ctex}
\usetikzlibrary{positioning}

\begin{document}

\begin{tikzpicture}[
    % 全局设置
    node distance = 1.6cm and 1.4cm,
    every node/.style = {
        draw,
        rounded corners,
        thick,
        align = center,
        font = \sffamily\small,
        minimum width = 2.0cm,
        minimum height = 0.9cm
    },
    % 样式定义
    centerStyle/.style = {
        fill = blue!45,
        text = white,
        font = \sffamily\bfseries\large,
        line width = 2pt,
        minimum size = 2.5cm,
        circle
    },
    branchTitle/.style = {
        font = \sffamily\bfseries,
        line width = 1.5pt
    },
    subNode/.style = {
        fill = white,
        font = \sffamily\small,
        minimum width = 2.0cm,
        minimum height = 0.8cm
    },
    leafNode/.style = {
        fill = white,
        font = \sffamily\footnotesize,
        minimum width = 1.8cm,
        minimum height = 0.7cm
    },
    arrowStyle/.style = {
        ->,
        thick,
        >=stealth,
        shorten >=2pt
    }
]

    % ================= 1. 中心节点 =================
    \node[centerStyle] (center) {权重衰减 \\ Weight Decay};

    % ================= 2. 上方分支：基本概念 (红) =================
    \node[branchTitle, fill=red!15, above=of center] (basic) {基本概念 \\ Basics};

    \node[subNode, above left=of basic] (l2) {平方和正则化 \\ Sum-of-Squares};
    \node[subNode, above=of basic] (grad) {梯度形式 \\ Gradient};
    \node[subNode, above right=of basic] (gauss) {高斯先验 \\ Gaussian Prior};
    \node[subNode, left=of basic, yshift=0cm] (eff) {有效参数量 \\ Effective Parameters};

    % ================= 3. 右侧分支：一致性正则化 (蓝) =================
    \node[branchTitle, fill=orange!15, right=of center] (consist) {一致性正则化 \\ Consistent Regularizers};

    \node[subNode, above=of consist] (trans) {线性变换 \\ Linear Transform};
    \node[subNode, above right=of consist] (scale) {权重重缩放 \\ Weight Rescaling};
    \node[subNode, right=of consist] (bias-shift) {偏置平移 \\ Bias Shift};
    \node[subNode, below right=of consist, yshift=0cm] (improper) {非正常先验 \\ Improper Prior};
    \node[subNode, below=of consist] (hyper) {四个超参数 \\ Four Hyperparameters};

    % ================= 4. 下方分支：广义权重衰减 (绿) =================
    \node[branchTitle, fill=green!15, below=of center] (general) {广义权重衰减 \\ Generalized Weight Decay};

    \node[subNode, below left=of general] (q) {幂次 $q$ \\ $\Omega(w)=\frac{\lambda}{2}\sum_{j=1}^{M}|w_j|^q$};
    \node[subNode, below=of general] (lasso) {Lasso \\ $q=1$};
    \node[subNode, below right=of general] (sparse) {稀疏性 \\ Sparsity};
    \node[subNode, left=of general, yshift=0cm] (constraint) {约束形式 \\ Constraint};

    % ================= 5. 左侧分支：模型复杂度 (紫) =================
    \node[branchTitle, fill=purple!15, left=of center] (complex) {模型复杂度 \\ Model Complexity};

    \node[subNode, above left=of complex] (overfit) {控制过拟合 \\ Over-fitting};
    \node[subNode, left=of complex] (lambda) {正则系数 $\lambda$ \\ Regularization};
    \node[subNode, below left=of complex] (shift) {复杂度转移 \\ Complexity Shift};
    
    % ================= 连线逻辑 =================
    % 中心 -> 四大分支标题
    \draw[arrowStyle, red] (center) -- (basic);
    \draw[arrowStyle, blue] (center) -- (consist);
    \draw[arrowStyle, green!60!black] (center) -- (general);
    \draw[arrowStyle, purple] (center) -- (complex);

    % 分支标题 -> 子节点 (上方)
    \draw[arrowStyle, red!60] (basic) -- (l2);
    \draw[arrowStyle, red!60] (basic) -- (grad);
    \draw[arrowStyle, red!60] (basic) -- (gauss);
    \draw[arrowStyle, red!60] (basic) -- (eff);

    % 分支标题 -> 子节点 (右侧)
    \draw[arrowStyle, blue!60] (consist) -- (trans);
    \draw[arrowStyle, blue!60] (consist) -- (scale);
    \draw[arrowStyle, blue!60] (consist) -- (bias-shift);
    \draw[arrowStyle, blue!60] (consist) -- (improper);
    \draw[arrowStyle, blue!60] (consist) -- (hyper);

    % 分支标题 -> 子节点 (下方)
    \draw[arrowStyle, green!60!black] (general) -- (q);
    \draw[arrowStyle, green!60!black] (general) -- (lasso);
    \draw[arrowStyle, green!60!black] (general) -- (sparse);
    \draw[arrowStyle, green!60!black] (general) -- (constraint);

    % 分支标题 -> 子节点 (左侧)
    \draw[arrowStyle, purple!60] (complex) -- (overfit);
    \draw[arrowStyle, purple!60] (complex) -- (lambda);
    \draw[arrowStyle, purple!60] (complex) -- (shift);
    
\end{tikzpicture}

\end{document}